% Independent mathematical audit of horn_candidate.py.
% Includable notes; parent report should rename labels if necessary.
\section{Independent audit of the unconditional-clause presentation}

\paragraph{Scope of the claim.}
The proposal has four explicit inference schemes and one fixed initial
clause. It has not eliminated reflexivity as a mathematical assumption:
the initial clause states universal reflexivity. The improvement is that
one uniform clause-use operation handles universal reflexivity, ordinary
hypotheses, and conditional closures, without a separate reflexivity
instruction or proofs of a neutral antecedent for ordinary hypotheses.

\subsection{Definition}
Use zero-normal iterator terms and formulas
\[
 t::=x\mid 0\mid St\mid I(n,a;u.h),\qquad
 F::=(s=t)\mid D_{x,y}F.
\]
A term is zero-normal when no iterator has literal-zero counter.
Every substitution is followed by zero normalization.
A clause has the form
\[
 [\,A_1,\ldots,A_k\,;\,y.B(y)\,],\qquad k\in\{0,1\},
\]
where the bound variable \(y\) is absent from every antecedent.
Its reading is
\[
 \bigwedge_{j=1}^k A_j^* \longrightarrow \forall y\,B(y)^*,
\]
with the empty-antecedent reading simply \(\forall y B(y)^*\).
The ordinary reading of \(D\) remains
\[
 (D_{x,y}B)^*=\exists x\bigl(B(x,0)^*\to\forall y B(x,y)^*\bigr).
\]
Every context contains the fixed clause
\[
 \mathcal R=[\,\,;\,y.(y=y)\,].
\]
An ordinary assumption \(A\) is the unconditional constant clause
\(\langle A\rangle=[\,\,;\,y.A\,]\), with \(y\) fresh for \(A\).

The four rules are:
\[
 \frac{[A_1,\ldots,A_k;y.B]\in\Gamma\quad
       \Gamma\vdash A_1\quad\cdots\quad\Gamma\vdash A_k}
      {\Gamma\vdash B\langle t/y\rangle}
 \quad(\mathrm{Use}),
\]
\[
 \frac{\Gamma\vdash I(c,a;u.h)=t\qquad
       \Gamma\vdash I(Sc,a;u.h)=r}
      {\Gamma\vdash r=h\langle t/u\rangle}
 \quad(\mathrm{Step}),
\]
\[
 \frac{\Gamma\vdash D_{x,y}B\qquad
       \Gamma,[B(b,0);y.B(b,y)]\vdash A}
      {\Gamma\vdash A}
 \quad(\mathrm{Open}),
\]
\[
 \frac{\Gamma,\langle B(t,\eta)\rangle\vdash B(t,Sz)}
      {\Gamma\vdash D_{x,y}B}
 \quad(\mathrm{Close}),
\qquad
 \eta=\begin{cases}
 z,& B\text{ is a literal equation and }t\text{ is literal }0,\\
 0,&\text{otherwise}.
 \end{cases}
\]
All instantiations in these displays include zero normalization.
The parameters \(b\) in Open and \(z\) in Close are fresh for the
outer context and outer conclusion, and \(z\) is also fresh for \(t\).
The \(\mathcal R\) clause is retained in every premise.

\subsection{Exact comparison with the five-pattern calculus}
The zero-normal five-pattern calculus has the same Step, Open and
syntax-directed Close, together with explicit Refl and conditional clause
use. Its ordinary assumption clause is
\([0=0;y.A]\).

Translate an unconditional new clause \([\,\,;y.B]\) to the old
conditional clause \([0=0;y.B]\); translate conditional clauses
identically. Translate the distinguished initial \(\mathcal R\)
operationally, rather than making it an extra old assumption: whenever
Use selects \(\mathcal R\), use old Refl at the supplied target.

For a new Use of another unconditional clause, old Refl proves \(0=0\),
so old clause use returns exactly the same normalized body instance.
Conditional uses are identical. Step and Open are unchanged, and a
new local unconditional assumption introduced by Close is translated to
the old neutral-antecedent assumption. This is a structural induction on
the new derivation, proving soundness with arbitrary retained side
clauses.

Conversely, old Refl at \(t\) becomes new Use of \(\mathcal R\) at \(t\),
with an empty premise list. An old local assumption
\([0=0;y.A]\) created by Close becomes the new unconditional assumption
\([\,\,;y.A]\). Uses of this distinguished occurrence discard the old
proof of \(0=0\). Every other old conditional clause and its Use is
translated unchanged. Step, Open and Close commute with this translation.
The local identities of context occurrences are retained, so there is
no ambiguity even when two clauses happen to be syntactically equal.

Thus the two calculi derive the same public consequences under translated
contexts. In particular, they prove exactly the same closed Eq/\(D\)
formulas. Both maps are effective finite proof transformations.
No general substitution theorem for syntax-directed Close has been used.

\subsection{Why zero normalization and the witness test are safe}
Write \(N\) for recursive contraction of literal-zero-counter iterators.
It preserves free variables and satisfies
\[
 N(Nt)=Nt,\qquad
 N(t[\sigma])=N(Nt[N\circ\sigma]).
\]
Consequently hereditary opening is associative at the level of normal
forms. The successor computation is obtained by Step on two
reflexivities, each supplied by Use of \(\mathcal R\).
With \(J_s=I(S0,s;u.u)\), two Steps give the pivot rule
\[
 s=t,\ s=r\ \vdash\ r=t.
\]
This supplies symmetry and transitivity. Applying Step to
\(I(S0,s;u.h)\) gives congruence in every term template and hence transport
in literal equations. These arithmetic derivations use no induction
or logical rule.

The nonliteral canonical zero
\[
 q=I(S0,0;u.u)
\]
satisfies \(q=0\), by one Step, while its syntax is not literal zero.
This repairs the exceptional atomic case introduction. Suppose
\[
 \Gamma,\langle E(0,0)\rangle\vdash E(0,Sz).
\]
First apply Close to the dummy-witness matrix \(E(0,y)\) with witness
\(q\). Since the witness is not literal zero, Close uses its case mode
and gives \(D_{x,y}E(0,y)\). Open it to retain the canonical clause
\([E(0,0);y.E(0,y)]\).
Now apply Close to \(D_{x,y}E(x,y)\), again at witness \(q\).
Its new assumption is \(E(q,0)\). Transport along \(q=0\) gives
\(E(0,0)\), clause use gives \(E(0,Sz)\), and transport along \(0=q\)
gives \(E(q,Sz)\). The fresh witness of the outer Open is absent from
the final conclusion. This derives the omitted atomic case rule at
literal witness zero.

For equation progress, a dummy-witness matrix \(E(y)\) and literal
witness zero directly give \(D_{x,y}E(y)\), the conditional-closure
abbreviation. Therefore the kernel's independent case and equation
progress rules are both available.

\subsection{Interpretations and consistency}
The new public language and its ordinary realization are identical to
those of the five-pattern zero-normal presentation. The new initial
clause is the ordinary equality theorem \(\forall y(y=y)\).
The exact proof equivalence above therefore transfers the established
native-to-ordinary equivalence and the ordinary scalar bi-interpretation
with iPRA. It changes neither the numerical domain nor equality.
The numerical comparison maps are the same identity-graph comparisons
as in the baseline; both term/formula round-trip theorems are inherited
after these explicit proof translations. Translating a derivation of
\(0=S0\) in either direction gives equiconsistency.

This is stronger than mere agreement in the standard natural numbers:
the proof transformations work with arbitrary retained assumptions and
do not assume standardness of numerical models.

\subsection{Audit of the Python scoping representation}
In \texttt{horn\_candidate.py}, negative integers are absolute free
levels, positive integers are indices of syntactically enclosing binders,
and zero is the arithmetic constant. A term supplied at scope \(s\)
has free levels in \(\{-s,\ldots,-1\}\) and no free positive index.
Opening inserts only such terms, so an inserted value needs no index
shift under intervening object-language binders. The traversal increases
its cutoff by one in an iterator body and by two in a \(D\)-matrix.
It removes the selected binder by decrementing larger positive indices.

Increasing scope from \(s\) to \(s+1\) introduces the level \(-s-1\);
old negative levels retain their meanings, so outer contexts need no
renaming. In Open, its target is validated at the old scope. In Close,
the displayed \(D\) goal and witness are validated at the old scope, and
the target is constructed before the branch is checked. These facts
prevent eigenvariable escape. The initial clause body \(1=1\) has one
bound variable and is well scoped.

Typing is preserved inductively by each successful inference. Step's
pattern matching obtains its templates only from two already-derived
well-formed equations; its literal-source check enforces a common
iterator counter, seed, and body. Use compares the complete ordered list
of derived antecedents to the stored list, so a missing or extra premise
cannot be accepted. Invalid JSON shapes or scoping failures raise an
exception and are rejected by the command-line wrapper.

\subsection{Limits}
The four rules plus initial clause are a smaller representation of the
same mathematics. They are not a proof that universal reflexivity is
derivable from Step alone, and they are not an absolute minimality
theorem. The normalization and initial clause remain trusted parts of
the complete specification and must be counted in source measurements.


\section{Exact flattening of the unary formula spine}
\subsection{The ranked formula grammar}
The public grammar has only an equality leaf and a unary \(D\) constructor.
Consequently every formula has the unique form \(D^k(s=t)\), for a natural
number \(k\). Each \(D\) binds its own ordered pair of variables. This is
an exact syntactic fact, without logical equivalence or normalization.
In de Bruijn notation encode it by
\[
 \operatorname{code}(D^k(s=t))=(k,s,t).
\]
The numerical term codes are unchanged. At external scope \(m\), this
triple is valid exactly when \(k\) is a nonnegative integer and its two
terms are valid at external scope \(m\) with \(2k\) surrounding bound
variables. Under a clause binder the depth is \(2k+1\).
The decoder is
\[
 \operatorname{decode}(0,s,t)=(s=t),\qquad
 \operatorname{decode}(k+1,s,t)=D(\operatorname{decode}(k,s,t)).
\]
Terms at the equality leaf refer to all enclosing binders, so no index
is shifted by decoding. Structural induction gives exact inverses:
\[
 \operatorname{code}(\operatorname{decode}(k,s,t))=(k,s,t),\qquad
 \operatorname{decode}(\operatorname{code}(A))=A.
\]
Starting with named variables, the second equation means equality after
the standard de Bruijn encoding, equivalently up to alpha-renaming.

\subsection{Opening and the order of the bound variables}
Let \(o_d(t,v)\) be hereditary term opening at cutoff \(d\): replace
index \(d+1\) by \(v\), decrement higher indices, account for iterator
binders, and contract newly exposed zero iterators. Supplied values have
no free positive index. Ranked formula opening is
\[
 O_d((k,s,t),v)
   =(k,o_{d+2k}(s,v),o_{d+2k}(t,v)).
\]
This exactly matches old formula opening, which increases its cutoff
by two at each \(D\) before reaching the equality leaf. Renaming and
simultaneous substitution satisfy the analogous componentwise lemma.

For a ranked \(D\) formula \((k,s,t)\), with \(k>0\), its matrix is
\(M=(k-1,s,t)\). The outer \(y\) has term index \(2k-1\), while the
outer \(x\) has index \(2k\). Therefore \(B_b=O_1(M,b)\) substitutes
for \(x\) and leaves \(y\) as the one external clause binder.
Its antecedent is \(O_0(B_b,0)\), and its instance at \(w\) is
\(O_0(B_b,w)\). In particular ranked Open adds
\[
 [\,O_0(B_b,0)\,;\,B_b\,],\qquad B_b=O_1((k-1,s,t),b).
\]
For Close at witness \(a\), put \(B_a=O_1((k-1,s,t),a)\).
The local unconditional assumption is \(O_0(B_a,\eta)\), and the
branch target is \(O_0(B_a,Sz)\). The matrix is a literal equation
exactly when \(k=1\). Thus the mode test becomes
\[
 \eta=z\ \Longleftrightarrow\ (k=1\text{ and }a=0\text{ literally}),
 \qquad \eta=0\text{ otherwise}.
\]
These formulas account for all binder ordering and index decrements.

\subsection{Proof equivalence and interpretation preservation}
Map every formula in a proof, clause, and goal to its ranked code;
leave numerical terms and proof-tree shape unchanged. The opening lemma
shows that Use, Open and Close commute exactly with this map.
The initial reflexivity clause becomes
\([\,\,;(0,1,1)]\).
Step is unchanged except that equation codes have rank zero.
The validity predicate is preserved by the \(2k\) depth criterion.
Induction on derivations therefore gives
\[
 \Gamma\vdash A\quad\Longleftrightarrow\quad
 \operatorname{code}(\Gamma)\vdash_{\rm ranked}\operatorname{code}(A).
\]
Componentwise decoding is the inverse proof map. This is a syntax
isomorphism, stronger than a mere consequence-equivalence theorem.

Define the ordinary interpretation of a triple as that of its decoder.
Then interpretation preservation is an identity:
\[
 (k,s,t)^*=(\operatorname{decode}(k,s,t))^*.
\]
All arithmetic derivations, the nonliteral-zero witness repair, and the
bi-interpretation with iPRA transfer verbatim. There is no new satisfaction
predicate, evaluator, choice principle, or induction instance.

\paragraph{Exact scope of the simplification.}
A recursive formula datatype is replaced by a natural rank and two terms.
Formula validation and opening are two term operations at a computed
depth. The arithmetic grammar and meaning of \(D\) are unchanged.
A binary-coded rank compactly represents a finite string of binders;
decoding can increase data length. This is a transparent certificate
representation choice, not an implicit theorem oracle. Universal
reflexivity remains in the fixed initial clause and is still counted.

